% =====================================================================
%  Week 2 Implementation — Adaptive Retrieval Gating (core novelty)
%  drop-in chapter section for the dissertation
%  Project: Adaptive Retrieval Gating for Safety-Aware Medical RAG
%  Author:  Tharit Mohamed Hussain Khan (22058691), KCL MSc AI, 7CCSMPRJ
%
%  To include in the main report:  % =====================================================================
%  Week 2 Implementation — Adaptive Retrieval Gating (core novelty)
%  drop-in chapter section for the dissertation
%  Project: Adaptive Retrieval Gating for Safety-Aware Medical RAG
%  Author:  Tharit Mohamed Hussain Khan (22058691), KCL MSc AI, 7CCSMPRJ
%
%  To include in the main report:  % =====================================================================
%  Week 2 Implementation — Adaptive Retrieval Gating (core novelty)
%  drop-in chapter section for the dissertation
%  Project: Adaptive Retrieval Gating for Safety-Aware Medical RAG
%  Author:  Tharit Mohamed Hussain Khan (22058691), KCL MSc AI, 7CCSMPRJ
%
%  To include in the main report:  % =====================================================================
%  Week 2 Implementation — Adaptive Retrieval Gating (core novelty)
%  drop-in chapter section for the dissertation
%  Project: Adaptive Retrieval Gating for Safety-Aware Medical RAG
%  Author:  Tharit Mohamed Hussain Khan (22058691), KCL MSc AI, 7CCSMPRJ
%
%  To include in the main report:  \input{week2_gating}
%  Required preamble packages (shared with week1_implementation.tex):
%    \usepackage{booktabs} \usepackage{amsmath} \usepackage{amssymb}
%
%  NOTE: the prototype results table (Table~\ref{tab:p5-prototype}) is a
%  PLACEHOLDER pending completion of the 200-question P5 run. Fill the marked
%  cells from results/raw_logs/p5_mirage200.jsonl once the run finishes.
% =====================================================================

\section{Adaptive Retrieval Gating (Week 2)}
\label{sec:impl-week2}

This section documents the project's core novel contribution: the
training-free retrieval gates that decide, per query, whether to retrieve
external evidence or answer from the model's parametric knowledge. Building on
the foundation layer of Section~\ref{sec:impl-week1}, this sprint implemented
three independent gating signals, combined them into a majority-vote ensemble,
and composed them into the selective-retrieval policy~(P5). The result is a
working prototype in which P5 runs end-to-end on the local quantised model and
is directly comparable to the closed-book floor established in Week~1.

% ---------------------------------------------------------------------
\subsection{Design Rationale: Why Three Gates}
\label{sec:gate-rationale}

A single uncertainty signal is fragile. Token entropy can be inflated by
benign lexical variation (synonyms, formatting tokens) on a query the model
actually knows; the logit margin can be large on a confidently \emph{wrong}
answer; self-reported confidence is poorly calibrated in small models. Each
signal therefore has a characteristic failure mode in which it would gate
incorrectly. The design response is an ensemble of three signals derived from
\emph{different} aspects of the model's behaviour, combined by majority vote so
that no single signal's failure mode dominates the decision. This mirrors the
project's safety-first framing: in a clinical setting, the cost of an
unnecessary retrieval (latency, energy) is far lower than the cost of a
confidently wrong closed-book answer, so the ensemble is deliberately biased
towards retrieving when the signals disagree.

% ---------------------------------------------------------------------
\subsection{Gate 1: Token Entropy}
\label{sec:gate-entropy}

The entropy gate generates a short draft answer (48 tokens) without retrieval
and measures the model's uncertainty from the full-vocabulary logit
distribution at each generated position. For generated token $t$ with softmax
distribution $p_t$ over the vocabulary $V$, the per-token Shannon entropy (in
nats) and the gate signal are
\begin{align}
  H(p_t) &= -\sum_{v \in V} p_t(v)\,\log p_t(v), &
  \bar{H} &= \frac{1}{N}\sum_{t=1}^{N} H(p_t),
\end{align}
and the gate retrieves when $\bar{H} > \tau_H$ (default $\tau_H = 2.5$). High
mean entropy indicates the model is spreading probability mass over many
candidate continuations --- a direct measure of lexical-level uncertainty.

Beyond the scalar $\bar{H}$ used for the decision, the gate also returns the
full per-token array $\{H(p_t)\}_{t=1}^N$. This array is the basis of the
\emph{Token-Level Entropy Attribution} contribution: it localises uncertainty
to specific tokens of the draft, revealing \emph{which part} of a multi-part
clinical question the model is unsure about rather than merely \emph{that} it
is unsure. The array is propagated through the policy result into the JSONL
audit log, where it is available to the planned interpretability
visualisation. The gate requires the raw logit buffer
($\texttt{logits\_all=True}$); when that buffer is unavailable on the low
hardware tier, the gate abstains rather than producing a spurious signal.

% ---------------------------------------------------------------------
\subsection{Gate 2: Logit Margin}
\label{sec:gate-margin}

The margin gate measures the model's \emph{decisiveness} from the gap between
its two most probable tokens at each position. For generated token $t$ with
top-two probabilities $p_t^{(1)} \ge p_t^{(2)}$,
\begin{align}
  m_t &= p_t^{(1)} - p_t^{(2)}, &
  \bar{M} &= \frac{1}{N}\sum_{t=1}^{N} m_t,
\end{align}
and the gate retrieves when $\bar{M} < \tau_M$ (default $\tau_M = 0.3$). A
small margin means the model is torn between two near-equiprobable
continuations, a complementary notion to entropy: whereas entropy summarises
the whole distribution, the margin focuses on the decision boundary between the
top two candidates.

Crucially, the margin gate obtains its signal from the completion API's
top-$k$ log-probabilities rather than the raw logit buffer, a lighter
dependency than the entropy gate's. The parser that converts the API's
per-token log-probability dictionaries into margins was validated against the
real inference library during this sprint, confirming that the live output
shape matches the mock used in testing, so that the unit tests are faithful to
production behaviour. Basis: the margin signal of TARG.

% ---------------------------------------------------------------------
\subsection{Gate 3: Hallucination Probe}
\label{sec:gate-probe}

The hallucination-probe gate replaces the verbalized-confidence gate of the
original design. Rather than asking the model to introspect on its own
confidence --- a task at which small models are unreliable --- it tests whether
the model's answer is \emph{stable} under reframing. Two drafts are generated
from the same question under two different prompt framings:
\begin{align*}
  \text{draft}_A &= \mathcal{M}\bigl(\text{probe}_A(q)\bigr), &
  \text{draft}_B &= \mathcal{M}\bigl(\text{probe}_B(q)\bigr),
\end{align*}
where $\text{probe}_A$ adopts an assistant-role framing and $\text{probe}_B$ a
direct-question framing. For multiple-choice questions the gate compares the
extracted option letters and retrieves when they differ; for open-ended
questions it computes the token-level $F_1$ between the two drafts and
retrieves when $F_1 < 0.7$. The intuition is that genuinely held knowledge is
invariant to surface framing, whereas a guessed answer is sensitive to it, so
disagreement across framings is evidence of instability.

This gate is purely text-based: it never reads the logit buffer, and is
therefore the only gate that remains available on the low hardware tier where
$\texttt{logits\_all=False}$. Its cost is two draft generations per query, an
overhead reported explicitly in the evaluation. Basis: self-consistency (Wang
et al., 2023), here reduced from majority voting over many samples to a binary
agreement test between two framings.

% ---------------------------------------------------------------------
\subsection{The Ensemble and Hardware-Aware Degradation}
\label{sec:gate-ensemble}

The three gates are combined by majority vote. Let $g_i \in
\{\textsc{retrieve}, \textsc{skip}\}$ be the decision of member $i$ among the
\emph{available} members $A \subseteq \{\text{entropy}, \text{margin},
\text{probe}\}$. The ensemble retrieves when
\[
  \bigl|\{\, i \in A : g_i = \textsc{retrieve} \,\}\bigr| \;\ge\; v_{\min},
  \qquad v_{\min} = 2 .
\]
Only members reporting themselves available cast a vote; a gate that abstains
(the entropy and margin gates on the low tier) is excluded from both the
numerator and the count of voters.

This availability accounting produces a principled \emph{degraded mode}. On the
low tier, only the probe gate is available, so the strict threshold
$v_{\min}=2$ is unreachable. Rather than failing, the ensemble falls back to
``retrieve if any available member votes to retrieve'' and records a
\texttt{degraded} flag in the audit log. The fallback is enforced at
construction time by the policy factory, which inspects the model's
$\texttt{logits\_all}$ setting and downgrades any requested gate to
probe-only when the logit buffer is disabled. The consequence is an important
design property: \emph{the gate a policy actually runs is constrained by the
hardware tier, not merely by the policy configuration file}. This makes the
hardware-aware behaviour auditable and turns the low-tier limitation into a
documented finding supporting the \emph{Hardware-Aware Policy Escalation}
contribution.

% ---------------------------------------------------------------------
\subsection{The Gated Policy (P5)}
\label{sec:p5-policy}

The selective-retrieval policy composes the ensemble with the retriever and the
language model. For each question it first obtains the gate decision; on
\textsc{retrieve} it behaves like the always-retrieve baseline (retrieve
evidence, answer with it in context), and on \textsc{skip} it behaves like the
closed-book baseline (answer from parametric knowledge alone). The complete
gate decision payload --- the scalar signal, the per-member votes, the
per-token entropy array, and the two probe drafts --- is attached to the policy
result and written into the per-query JSONL record, so that every gating
decision is fully reconstructable after the fact. Policy construction is
centralised in a factory that assembles the correct policy and gate from a
validated configuration, replacing the hand-wired construction used in Week~1.

% ---------------------------------------------------------------------
\subsection{Verification}
\label{sec:gate-verification}

Every gate and the gated policy were developed test-first against the mock
backend, whose synthetic peaked and uniform logit distributions exercise both
the \textsc{skip} and \textsc{retrieve} paths deterministically. The gate test
suite covers each signal's threshold behaviour, the per-token entropy
recording, the abstain paths under disabled logits, the ensemble vote
arithmetic, and the degraded-mode fallback; the policy test suite covers both
branches of P5, the propagation of gate fields into the result, and the
factory's construction of the correct gate per hardware tier (including the
low-tier probe-only downgrade). The full unit suite stands at 86 passing tests
and continues to execute without any model download, preserving the project's
fast, reproducible verification story.

% ---------------------------------------------------------------------
\subsection{Prototype Result}
\label{sec:p5-prototype}

The prototype was exercised by running P5 over the same 200 MIRAGE questions
used for the Week~1 closed-book baseline, using the local Llama-3.2-3B model
and a curated pilot retrieval corpus. Using the identical question set makes the
gated policy directly comparable to the closed-book floor of
Table~\ref{tab:p3-baseline}. The pilot corpus is a deliberately small,
curated set; retrieval-\emph{quality} figures obtained against it are
provisional and will be re-measured against the full corpus, but the
retrieval-\emph{rate}, gate-agreement and latency-overhead figures reported
here characterise the gating mechanism itself and are not sensitive to corpus
size. The aggregate prototype figures are given in
Table~\ref{tab:p5-prototype}, and the per-gate behaviour --- which proves the
most informative result of this sprint --- in Table~\ref{tab:p5-gatevotes}.

\begin{table}[ht]
  \centering
  \caption{Prototype P5 (3-gate ensemble, default thresholds) vs.\ the P3
           closed-book floor on the same 200 MIRAGE questions, Llama-3.2-3B
           Q4\_K\_M, medium tier. At the literature-default thresholds the gate
           never fired; the two policies are therefore identical in accuracy.
           See Table~\ref{tab:p5-gatevotes} for the cause.}
  \label{tab:p5-prototype}
  \begin{tabular}{lrr}
    \toprule
    \textbf{Metric} & \textbf{P3 (closed-book)} & \textbf{P5 (gated, uncal.)} \\
    \midrule
    Accuracy            & 62.0\% & 62.0\% \\
    Retrieval rate      & 0\%    & 0\% \\
    Median latency      & 7.42\,s & 143.7\,s \\
    95th-pct.\ latency  & 9.25\,s & 154.1\,s \\
    Mean energy / query & $6.77\times10^{-4}$\,kWh & $1.20\times10^{-3}$\,kWh \\
    \bottomrule
  \end{tabular}
\end{table}

\begin{table}[ht]
  \centering
  \caption{Per-gate signal statistics and vote counts over the 200 questions.
           The entropy and margin signals never cross their default thresholds
           on this quantised model, so only the hallucination probe ever votes
           to retrieve --- and a lone probe vote cannot reach the majority of
           two. This is a threshold-transfer finding, not a defect.}
  \label{tab:p5-gatevotes}
  \begin{tabular}{lrrrr}
    \toprule
    \textbf{Gate} & \textbf{Signal range} & \textbf{Default $\tau$} &
    \textbf{Retrieve votes} & \textbf{Skip votes} \\
    \midrule
    Entropy (nats)        & 0.27--1.59 & $>2.5$  & 0   & 200 \\
    Margin                & 0.45--0.89 & $<0.3$  & 0   & 200 \\
    Hallucination probe   & ---        & ---     & 80  & 120 \\
    \bottomrule
  \end{tabular}
\end{table}

\paragraph{Interpretation and the central finding.}
At the thresholds inherited from the general-domain literature ($\tau_H=2.5$
nats, $\tau_M=0.3$), the gate retrieved on \emph{zero} of the 200 questions, so
P5 reduced to the closed-book policy and reproduced its 62.0\% accuracy exactly.
The per-gate trace in Table~\ref{tab:p5-gatevotes} explains why: on the
4-bit-quantised Llama-3.2-3B, greedy drafts are sharply peaked, so the mean
token entropy (observed range 0.27--1.59 nats) never approaches the 2.5-nat
trigger, and the top-two probability margin (0.45--0.89) never falls below the
0.3 trigger. Only the hallucination probe --- which is threshold-free and
measures cross-framing consistency --- ever signalled uncertainty, voting to
retrieve on 80 of 200 questions; but under a majority rule requiring two of
three votes, a single dissenting gate can never carry the decision.

This is the most important result of the sprint, and it is a genuine empirical
finding rather than an implementation fault: \emph{gate thresholds calibrated
on general-domain models do not transfer to a small quantised medical model.}
It directly motivates the threshold-calibration sweep (Sprint~4): the observed
signal distributions suggest model-specific operating points near the
$75^{\text{th}}$ percentile of entropy ($\approx 0.9$ nats) and the
$25^{\text{th}}$ percentile of margin ($\approx 0.64$), which would move the
logit-based gates off their floor and let the three signals genuinely compete.
The prototype has thus done exactly what a prototype should: it validated the
end-to-end mechanism (gates compute, votes aggregate, the policy branches and
logs correctly) while surfacing a calibration requirement that becomes a
substantive point of analysis in the evaluation.

\paragraph{Latency overhead.}
The gated policy's median per-query latency of 143.7\,s against the
closed-book 7.42\,s reflects the cost of the gating drafts: the ensemble issues
the entropy/margin draft plus the two hallucination-probe drafts before the
final answer, and on CPU these dominate the per-query time. This overhead is
reported honestly as the price of the current ensemble and is itself an
argument for the calibration work and for the low-tier probe-only
configuration, where fewer gate calls are made.

%  Required preamble packages (shared with week1_implementation.tex):
%    \usepackage{booktabs} \usepackage{amsmath} \usepackage{amssymb}
%
%  NOTE: the prototype results table (Table~\ref{tab:p5-prototype}) is a
%  PLACEHOLDER pending completion of the 200-question P5 run. Fill the marked
%  cells from results/raw_logs/p5_mirage200.jsonl once the run finishes.
% =====================================================================

\section{Adaptive Retrieval Gating (Week 2)}
\label{sec:impl-week2}

This section documents the project's core novel contribution: the
training-free retrieval gates that decide, per query, whether to retrieve
external evidence or answer from the model's parametric knowledge. Building on
the foundation layer of Section~\ref{sec:impl-week1}, this sprint implemented
three independent gating signals, combined them into a majority-vote ensemble,
and composed them into the selective-retrieval policy~(P5). The result is a
working prototype in which P5 runs end-to-end on the local quantised model and
is directly comparable to the closed-book floor established in Week~1.

% ---------------------------------------------------------------------
\subsection{Design Rationale: Why Three Gates}
\label{sec:gate-rationale}

A single uncertainty signal is fragile. Token entropy can be inflated by
benign lexical variation (synonyms, formatting tokens) on a query the model
actually knows; the logit margin can be large on a confidently \emph{wrong}
answer; self-reported confidence is poorly calibrated in small models. Each
signal therefore has a characteristic failure mode in which it would gate
incorrectly. The design response is an ensemble of three signals derived from
\emph{different} aspects of the model's behaviour, combined by majority vote so
that no single signal's failure mode dominates the decision. This mirrors the
project's safety-first framing: in a clinical setting, the cost of an
unnecessary retrieval (latency, energy) is far lower than the cost of a
confidently wrong closed-book answer, so the ensemble is deliberately biased
towards retrieving when the signals disagree.

% ---------------------------------------------------------------------
\subsection{Gate 1: Token Entropy}
\label{sec:gate-entropy}

The entropy gate generates a short draft answer (48 tokens) without retrieval
and measures the model's uncertainty from the full-vocabulary logit
distribution at each generated position. For generated token $t$ with softmax
distribution $p_t$ over the vocabulary $V$, the per-token Shannon entropy (in
nats) and the gate signal are
\begin{align}
  H(p_t) &= -\sum_{v \in V} p_t(v)\,\log p_t(v), &
  \bar{H} &= \frac{1}{N}\sum_{t=1}^{N} H(p_t),
\end{align}
and the gate retrieves when $\bar{H} > \tau_H$ (default $\tau_H = 2.5$). High
mean entropy indicates the model is spreading probability mass over many
candidate continuations --- a direct measure of lexical-level uncertainty.

Beyond the scalar $\bar{H}$ used for the decision, the gate also returns the
full per-token array $\{H(p_t)\}_{t=1}^N$. This array is the basis of the
\emph{Token-Level Entropy Attribution} contribution: it localises uncertainty
to specific tokens of the draft, revealing \emph{which part} of a multi-part
clinical question the model is unsure about rather than merely \emph{that} it
is unsure. The array is propagated through the policy result into the JSONL
audit log, where it is available to the planned interpretability
visualisation. The gate requires the raw logit buffer
($\texttt{logits\_all=True}$); when that buffer is unavailable on the low
hardware tier, the gate abstains rather than producing a spurious signal.

% ---------------------------------------------------------------------
\subsection{Gate 2: Logit Margin}
\label{sec:gate-margin}

The margin gate measures the model's \emph{decisiveness} from the gap between
its two most probable tokens at each position. For generated token $t$ with
top-two probabilities $p_t^{(1)} \ge p_t^{(2)}$,
\begin{align}
  m_t &= p_t^{(1)} - p_t^{(2)}, &
  \bar{M} &= \frac{1}{N}\sum_{t=1}^{N} m_t,
\end{align}
and the gate retrieves when $\bar{M} < \tau_M$ (default $\tau_M = 0.3$). A
small margin means the model is torn between two near-equiprobable
continuations, a complementary notion to entropy: whereas entropy summarises
the whole distribution, the margin focuses on the decision boundary between the
top two candidates.

Crucially, the margin gate obtains its signal from the completion API's
top-$k$ log-probabilities rather than the raw logit buffer, a lighter
dependency than the entropy gate's. The parser that converts the API's
per-token log-probability dictionaries into margins was validated against the
real inference library during this sprint, confirming that the live output
shape matches the mock used in testing, so that the unit tests are faithful to
production behaviour. Basis: the margin signal of TARG.

% ---------------------------------------------------------------------
\subsection{Gate 3: Hallucination Probe}
\label{sec:gate-probe}

The hallucination-probe gate replaces the verbalized-confidence gate of the
original design. Rather than asking the model to introspect on its own
confidence --- a task at which small models are unreliable --- it tests whether
the model's answer is \emph{stable} under reframing. Two drafts are generated
from the same question under two different prompt framings:
\begin{align*}
  \text{draft}_A &= \mathcal{M}\bigl(\text{probe}_A(q)\bigr), &
  \text{draft}_B &= \mathcal{M}\bigl(\text{probe}_B(q)\bigr),
\end{align*}
where $\text{probe}_A$ adopts an assistant-role framing and $\text{probe}_B$ a
direct-question framing. For multiple-choice questions the gate compares the
extracted option letters and retrieves when they differ; for open-ended
questions it computes the token-level $F_1$ between the two drafts and
retrieves when $F_1 < 0.7$. The intuition is that genuinely held knowledge is
invariant to surface framing, whereas a guessed answer is sensitive to it, so
disagreement across framings is evidence of instability.

This gate is purely text-based: it never reads the logit buffer, and is
therefore the only gate that remains available on the low hardware tier where
$\texttt{logits\_all=False}$. Its cost is two draft generations per query, an
overhead reported explicitly in the evaluation. Basis: self-consistency (Wang
et al., 2023), here reduced from majority voting over many samples to a binary
agreement test between two framings.

% ---------------------------------------------------------------------
\subsection{The Ensemble and Hardware-Aware Degradation}
\label{sec:gate-ensemble}

The three gates are combined by majority vote. Let $g_i \in
\{\textsc{retrieve}, \textsc{skip}\}$ be the decision of member $i$ among the
\emph{available} members $A \subseteq \{\text{entropy}, \text{margin},
\text{probe}\}$. The ensemble retrieves when
\[
  \bigl|\{\, i \in A : g_i = \textsc{retrieve} \,\}\bigr| \;\ge\; v_{\min},
  \qquad v_{\min} = 2 .
\]
Only members reporting themselves available cast a vote; a gate that abstains
(the entropy and margin gates on the low tier) is excluded from both the
numerator and the count of voters.

This availability accounting produces a principled \emph{degraded mode}. On the
low tier, only the probe gate is available, so the strict threshold
$v_{\min}=2$ is unreachable. Rather than failing, the ensemble falls back to
``retrieve if any available member votes to retrieve'' and records a
\texttt{degraded} flag in the audit log. The fallback is enforced at
construction time by the policy factory, which inspects the model's
$\texttt{logits\_all}$ setting and downgrades any requested gate to
probe-only when the logit buffer is disabled. The consequence is an important
design property: \emph{the gate a policy actually runs is constrained by the
hardware tier, not merely by the policy configuration file}. This makes the
hardware-aware behaviour auditable and turns the low-tier limitation into a
documented finding supporting the \emph{Hardware-Aware Policy Escalation}
contribution.

% ---------------------------------------------------------------------
\subsection{The Gated Policy (P5)}
\label{sec:p5-policy}

The selective-retrieval policy composes the ensemble with the retriever and the
language model. For each question it first obtains the gate decision; on
\textsc{retrieve} it behaves like the always-retrieve baseline (retrieve
evidence, answer with it in context), and on \textsc{skip} it behaves like the
closed-book baseline (answer from parametric knowledge alone). The complete
gate decision payload --- the scalar signal, the per-member votes, the
per-token entropy array, and the two probe drafts --- is attached to the policy
result and written into the per-query JSONL record, so that every gating
decision is fully reconstructable after the fact. Policy construction is
centralised in a factory that assembles the correct policy and gate from a
validated configuration, replacing the hand-wired construction used in Week~1.

% ---------------------------------------------------------------------
\subsection{Verification}
\label{sec:gate-verification}

Every gate and the gated policy were developed test-first against the mock
backend, whose synthetic peaked and uniform logit distributions exercise both
the \textsc{skip} and \textsc{retrieve} paths deterministically. The gate test
suite covers each signal's threshold behaviour, the per-token entropy
recording, the abstain paths under disabled logits, the ensemble vote
arithmetic, and the degraded-mode fallback; the policy test suite covers both
branches of P5, the propagation of gate fields into the result, and the
factory's construction of the correct gate per hardware tier (including the
low-tier probe-only downgrade). The full unit suite stands at 86 passing tests
and continues to execute without any model download, preserving the project's
fast, reproducible verification story.

% ---------------------------------------------------------------------
\subsection{Prototype Result}
\label{sec:p5-prototype}

The prototype was exercised by running P5 over the same 200 MIRAGE questions
used for the Week~1 closed-book baseline, using the local Llama-3.2-3B model
and a curated pilot retrieval corpus. Using the identical question set makes the
gated policy directly comparable to the closed-book floor of
Table~\ref{tab:p3-baseline}. The pilot corpus is a deliberately small,
curated set; retrieval-\emph{quality} figures obtained against it are
provisional and will be re-measured against the full corpus, but the
retrieval-\emph{rate}, gate-agreement and latency-overhead figures reported
here characterise the gating mechanism itself and are not sensitive to corpus
size. The aggregate prototype figures are given in
Table~\ref{tab:p5-prototype}, and the per-gate behaviour --- which proves the
most informative result of this sprint --- in Table~\ref{tab:p5-gatevotes}.

\begin{table}[ht]
  \centering
  \caption{Prototype P5 (3-gate ensemble, default thresholds) vs.\ the P3
           closed-book floor on the same 200 MIRAGE questions, Llama-3.2-3B
           Q4\_K\_M, medium tier. At the literature-default thresholds the gate
           never fired; the two policies are therefore identical in accuracy.
           See Table~\ref{tab:p5-gatevotes} for the cause.}
  \label{tab:p5-prototype}
  \begin{tabular}{lrr}
    \toprule
    \textbf{Metric} & \textbf{P3 (closed-book)} & \textbf{P5 (gated, uncal.)} \\
    \midrule
    Accuracy            & 62.0\% & 62.0\% \\
    Retrieval rate      & 0\%    & 0\% \\
    Median latency      & 7.42\,s & 143.7\,s \\
    95th-pct.\ latency  & 9.25\,s & 154.1\,s \\
    Mean energy / query & $6.77\times10^{-4}$\,kWh & $1.20\times10^{-3}$\,kWh \\
    \bottomrule
  \end{tabular}
\end{table}

\begin{table}[ht]
  \centering
  \caption{Per-gate signal statistics and vote counts over the 200 questions.
           The entropy and margin signals never cross their default thresholds
           on this quantised model, so only the hallucination probe ever votes
           to retrieve --- and a lone probe vote cannot reach the majority of
           two. This is a threshold-transfer finding, not a defect.}
  \label{tab:p5-gatevotes}
  \begin{tabular}{lrrrr}
    \toprule
    \textbf{Gate} & \textbf{Signal range} & \textbf{Default $\tau$} &
    \textbf{Retrieve votes} & \textbf{Skip votes} \\
    \midrule
    Entropy (nats)        & 0.27--1.59 & $>2.5$  & 0   & 200 \\
    Margin                & 0.45--0.89 & $<0.3$  & 0   & 200 \\
    Hallucination probe   & ---        & ---     & 80  & 120 \\
    \bottomrule
  \end{tabular}
\end{table}

\paragraph{Interpretation and the central finding.}
At the thresholds inherited from the general-domain literature ($\tau_H=2.5$
nats, $\tau_M=0.3$), the gate retrieved on \emph{zero} of the 200 questions, so
P5 reduced to the closed-book policy and reproduced its 62.0\% accuracy exactly.
The per-gate trace in Table~\ref{tab:p5-gatevotes} explains why: on the
4-bit-quantised Llama-3.2-3B, greedy drafts are sharply peaked, so the mean
token entropy (observed range 0.27--1.59 nats) never approaches the 2.5-nat
trigger, and the top-two probability margin (0.45--0.89) never falls below the
0.3 trigger. Only the hallucination probe --- which is threshold-free and
measures cross-framing consistency --- ever signalled uncertainty, voting to
retrieve on 80 of 200 questions; but under a majority rule requiring two of
three votes, a single dissenting gate can never carry the decision.

This is the most important result of the sprint, and it is a genuine empirical
finding rather than an implementation fault: \emph{gate thresholds calibrated
on general-domain models do not transfer to a small quantised medical model.}
It directly motivates the threshold-calibration sweep (Sprint~4): the observed
signal distributions suggest model-specific operating points near the
$75^{\text{th}}$ percentile of entropy ($\approx 0.9$ nats) and the
$25^{\text{th}}$ percentile of margin ($\approx 0.64$), which would move the
logit-based gates off their floor and let the three signals genuinely compete.
The prototype has thus done exactly what a prototype should: it validated the
end-to-end mechanism (gates compute, votes aggregate, the policy branches and
logs correctly) while surfacing a calibration requirement that becomes a
substantive point of analysis in the evaluation.

\paragraph{Latency overhead.}
The gated policy's median per-query latency of 143.7\,s against the
closed-book 7.42\,s reflects the cost of the gating drafts: the ensemble issues
the entropy/margin draft plus the two hallucination-probe drafts before the
final answer, and on CPU these dominate the per-query time. This overhead is
reported honestly as the price of the current ensemble and is itself an
argument for the calibration work and for the low-tier probe-only
configuration, where fewer gate calls are made.

%  Required preamble packages (shared with week1_implementation.tex):
%    \usepackage{booktabs} \usepackage{amsmath} \usepackage{amssymb}
%
%  NOTE: the prototype results table (Table~\ref{tab:p5-prototype}) is a
%  PLACEHOLDER pending completion of the 200-question P5 run. Fill the marked
%  cells from results/raw_logs/p5_mirage200.jsonl once the run finishes.
% =====================================================================

\section{Adaptive Retrieval Gating (Week 2)}
\label{sec:impl-week2}

This section documents the project's core novel contribution: the
training-free retrieval gates that decide, per query, whether to retrieve
external evidence or answer from the model's parametric knowledge. Building on
the foundation layer of Section~\ref{sec:impl-week1}, this sprint implemented
three independent gating signals, combined them into a majority-vote ensemble,
and composed them into the selective-retrieval policy~(P5). The result is a
working prototype in which P5 runs end-to-end on the local quantised model and
is directly comparable to the closed-book floor established in Week~1.

% ---------------------------------------------------------------------
\subsection{Design Rationale: Why Three Gates}
\label{sec:gate-rationale}

A single uncertainty signal is fragile. Token entropy can be inflated by
benign lexical variation (synonyms, formatting tokens) on a query the model
actually knows; the logit margin can be large on a confidently \emph{wrong}
answer; self-reported confidence is poorly calibrated in small models. Each
signal therefore has a characteristic failure mode in which it would gate
incorrectly. The design response is an ensemble of three signals derived from
\emph{different} aspects of the model's behaviour, combined by majority vote so
that no single signal's failure mode dominates the decision. This mirrors the
project's safety-first framing: in a clinical setting, the cost of an
unnecessary retrieval (latency, energy) is far lower than the cost of a
confidently wrong closed-book answer, so the ensemble is deliberately biased
towards retrieving when the signals disagree.

% ---------------------------------------------------------------------
\subsection{Gate 1: Token Entropy}
\label{sec:gate-entropy}

The entropy gate generates a short draft answer (48 tokens) without retrieval
and measures the model's uncertainty from the full-vocabulary logit
distribution at each generated position. For generated token $t$ with softmax
distribution $p_t$ over the vocabulary $V$, the per-token Shannon entropy (in
nats) and the gate signal are
\begin{align}
  H(p_t) &= -\sum_{v \in V} p_t(v)\,\log p_t(v), &
  \bar{H} &= \frac{1}{N}\sum_{t=1}^{N} H(p_t),
\end{align}
and the gate retrieves when $\bar{H} > \tau_H$ (default $\tau_H = 2.5$). High
mean entropy indicates the model is spreading probability mass over many
candidate continuations --- a direct measure of lexical-level uncertainty.

Beyond the scalar $\bar{H}$ used for the decision, the gate also returns the
full per-token array $\{H(p_t)\}_{t=1}^N$. This array is the basis of the
\emph{Token-Level Entropy Attribution} contribution: it localises uncertainty
to specific tokens of the draft, revealing \emph{which part} of a multi-part
clinical question the model is unsure about rather than merely \emph{that} it
is unsure. The array is propagated through the policy result into the JSONL
audit log, where it is available to the planned interpretability
visualisation. The gate requires the raw logit buffer
($\texttt{logits\_all=True}$); when that buffer is unavailable on the low
hardware tier, the gate abstains rather than producing a spurious signal.

% ---------------------------------------------------------------------
\subsection{Gate 2: Logit Margin}
\label{sec:gate-margin}

The margin gate measures the model's \emph{decisiveness} from the gap between
its two most probable tokens at each position. For generated token $t$ with
top-two probabilities $p_t^{(1)} \ge p_t^{(2)}$,
\begin{align}
  m_t &= p_t^{(1)} - p_t^{(2)}, &
  \bar{M} &= \frac{1}{N}\sum_{t=1}^{N} m_t,
\end{align}
and the gate retrieves when $\bar{M} < \tau_M$ (default $\tau_M = 0.3$). A
small margin means the model is torn between two near-equiprobable
continuations, a complementary notion to entropy: whereas entropy summarises
the whole distribution, the margin focuses on the decision boundary between the
top two candidates.

Crucially, the margin gate obtains its signal from the completion API's
top-$k$ log-probabilities rather than the raw logit buffer, a lighter
dependency than the entropy gate's. The parser that converts the API's
per-token log-probability dictionaries into margins was validated against the
real inference library during this sprint, confirming that the live output
shape matches the mock used in testing, so that the unit tests are faithful to
production behaviour. Basis: the margin signal of TARG.

% ---------------------------------------------------------------------
\subsection{Gate 3: Hallucination Probe}
\label{sec:gate-probe}

The hallucination-probe gate replaces the verbalized-confidence gate of the
original design. Rather than asking the model to introspect on its own
confidence --- a task at which small models are unreliable --- it tests whether
the model's answer is \emph{stable} under reframing. Two drafts are generated
from the same question under two different prompt framings:
\begin{align*}
  \text{draft}_A &= \mathcal{M}\bigl(\text{probe}_A(q)\bigr), &
  \text{draft}_B &= \mathcal{M}\bigl(\text{probe}_B(q)\bigr),
\end{align*}
where $\text{probe}_A$ adopts an assistant-role framing and $\text{probe}_B$ a
direct-question framing. For multiple-choice questions the gate compares the
extracted option letters and retrieves when they differ; for open-ended
questions it computes the token-level $F_1$ between the two drafts and
retrieves when $F_1 < 0.7$. The intuition is that genuinely held knowledge is
invariant to surface framing, whereas a guessed answer is sensitive to it, so
disagreement across framings is evidence of instability.

This gate is purely text-based: it never reads the logit buffer, and is
therefore the only gate that remains available on the low hardware tier where
$\texttt{logits\_all=False}$. Its cost is two draft generations per query, an
overhead reported explicitly in the evaluation. Basis: self-consistency (Wang
et al., 2023), here reduced from majority voting over many samples to a binary
agreement test between two framings.

% ---------------------------------------------------------------------
\subsection{The Ensemble and Hardware-Aware Degradation}
\label{sec:gate-ensemble}

The three gates are combined by majority vote. Let $g_i \in
\{\textsc{retrieve}, \textsc{skip}\}$ be the decision of member $i$ among the
\emph{available} members $A \subseteq \{\text{entropy}, \text{margin},
\text{probe}\}$. The ensemble retrieves when
\[
  \bigl|\{\, i \in A : g_i = \textsc{retrieve} \,\}\bigr| \;\ge\; v_{\min},
  \qquad v_{\min} = 2 .
\]
Only members reporting themselves available cast a vote; a gate that abstains
(the entropy and margin gates on the low tier) is excluded from both the
numerator and the count of voters.

This availability accounting produces a principled \emph{degraded mode}. On the
low tier, only the probe gate is available, so the strict threshold
$v_{\min}=2$ is unreachable. Rather than failing, the ensemble falls back to
``retrieve if any available member votes to retrieve'' and records a
\texttt{degraded} flag in the audit log. The fallback is enforced at
construction time by the policy factory, which inspects the model's
$\texttt{logits\_all}$ setting and downgrades any requested gate to
probe-only when the logit buffer is disabled. The consequence is an important
design property: \emph{the gate a policy actually runs is constrained by the
hardware tier, not merely by the policy configuration file}. This makes the
hardware-aware behaviour auditable and turns the low-tier limitation into a
documented finding supporting the \emph{Hardware-Aware Policy Escalation}
contribution.

% ---------------------------------------------------------------------
\subsection{The Gated Policy (P5)}
\label{sec:p5-policy}

The selective-retrieval policy composes the ensemble with the retriever and the
language model. For each question it first obtains the gate decision; on
\textsc{retrieve} it behaves like the always-retrieve baseline (retrieve
evidence, answer with it in context), and on \textsc{skip} it behaves like the
closed-book baseline (answer from parametric knowledge alone). The complete
gate decision payload --- the scalar signal, the per-member votes, the
per-token entropy array, and the two probe drafts --- is attached to the policy
result and written into the per-query JSONL record, so that every gating
decision is fully reconstructable after the fact. Policy construction is
centralised in a factory that assembles the correct policy and gate from a
validated configuration, replacing the hand-wired construction used in Week~1.

% ---------------------------------------------------------------------
\subsection{Verification}
\label{sec:gate-verification}

Every gate and the gated policy were developed test-first against the mock
backend, whose synthetic peaked and uniform logit distributions exercise both
the \textsc{skip} and \textsc{retrieve} paths deterministically. The gate test
suite covers each signal's threshold behaviour, the per-token entropy
recording, the abstain paths under disabled logits, the ensemble vote
arithmetic, and the degraded-mode fallback; the policy test suite covers both
branches of P5, the propagation of gate fields into the result, and the
factory's construction of the correct gate per hardware tier (including the
low-tier probe-only downgrade). The full unit suite stands at 86 passing tests
and continues to execute without any model download, preserving the project's
fast, reproducible verification story.

% ---------------------------------------------------------------------
\subsection{Prototype Result}
\label{sec:p5-prototype}

The prototype was exercised by running P5 over the same 200 MIRAGE questions
used for the Week~1 closed-book baseline, using the local Llama-3.2-3B model
and a curated pilot retrieval corpus. Using the identical question set makes the
gated policy directly comparable to the closed-book floor of
Table~\ref{tab:p3-baseline}. The pilot corpus is a deliberately small,
curated set; retrieval-\emph{quality} figures obtained against it are
provisional and will be re-measured against the full corpus, but the
retrieval-\emph{rate}, gate-agreement and latency-overhead figures reported
here characterise the gating mechanism itself and are not sensitive to corpus
size. The aggregate prototype figures are given in
Table~\ref{tab:p5-prototype}, and the per-gate behaviour --- which proves the
most informative result of this sprint --- in Table~\ref{tab:p5-gatevotes}.

\begin{table}[ht]
  \centering
  \caption{Prototype P5 (3-gate ensemble, default thresholds) vs.\ the P3
           closed-book floor on the same 200 MIRAGE questions, Llama-3.2-3B
           Q4\_K\_M, medium tier. At the literature-default thresholds the gate
           never fired; the two policies are therefore identical in accuracy.
           See Table~\ref{tab:p5-gatevotes} for the cause.}
  \label{tab:p5-prototype}
  \begin{tabular}{lrr}
    \toprule
    \textbf{Metric} & \textbf{P3 (closed-book)} & \textbf{P5 (gated, uncal.)} \\
    \midrule
    Accuracy            & 62.0\% & 62.0\% \\
    Retrieval rate      & 0\%    & 0\% \\
    Median latency      & 7.42\,s & 143.7\,s \\
    95th-pct.\ latency  & 9.25\,s & 154.1\,s \\
    Mean energy / query & $6.77\times10^{-4}$\,kWh & $1.20\times10^{-3}$\,kWh \\
    \bottomrule
  \end{tabular}
\end{table}

\begin{table}[ht]
  \centering
  \caption{Per-gate signal statistics and vote counts over the 200 questions.
           The entropy and margin signals never cross their default thresholds
           on this quantised model, so only the hallucination probe ever votes
           to retrieve --- and a lone probe vote cannot reach the majority of
           two. This is a threshold-transfer finding, not a defect.}
  \label{tab:p5-gatevotes}
  \begin{tabular}{lrrrr}
    \toprule
    \textbf{Gate} & \textbf{Signal range} & \textbf{Default $\tau$} &
    \textbf{Retrieve votes} & \textbf{Skip votes} \\
    \midrule
    Entropy (nats)        & 0.27--1.59 & $>2.5$  & 0   & 200 \\
    Margin                & 0.45--0.89 & $<0.3$  & 0   & 200 \\
    Hallucination probe   & ---        & ---     & 80  & 120 \\
    \bottomrule
  \end{tabular}
\end{table}

\paragraph{Interpretation and the central finding.}
At the thresholds inherited from the general-domain literature ($\tau_H=2.5$
nats, $\tau_M=0.3$), the gate retrieved on \emph{zero} of the 200 questions, so
P5 reduced to the closed-book policy and reproduced its 62.0\% accuracy exactly.
The per-gate trace in Table~\ref{tab:p5-gatevotes} explains why: on the
4-bit-quantised Llama-3.2-3B, greedy drafts are sharply peaked, so the mean
token entropy (observed range 0.27--1.59 nats) never approaches the 2.5-nat
trigger, and the top-two probability margin (0.45--0.89) never falls below the
0.3 trigger. Only the hallucination probe --- which is threshold-free and
measures cross-framing consistency --- ever signalled uncertainty, voting to
retrieve on 80 of 200 questions; but under a majority rule requiring two of
three votes, a single dissenting gate can never carry the decision.

This is the most important result of the sprint, and it is a genuine empirical
finding rather than an implementation fault: \emph{gate thresholds calibrated
on general-domain models do not transfer to a small quantised medical model.}
It directly motivates the threshold-calibration sweep (Sprint~4): the observed
signal distributions suggest model-specific operating points near the
$75^{\text{th}}$ percentile of entropy ($\approx 0.9$ nats) and the
$25^{\text{th}}$ percentile of margin ($\approx 0.64$), which would move the
logit-based gates off their floor and let the three signals genuinely compete.
The prototype has thus done exactly what a prototype should: it validated the
end-to-end mechanism (gates compute, votes aggregate, the policy branches and
logs correctly) while surfacing a calibration requirement that becomes a
substantive point of analysis in the evaluation.

\paragraph{Latency overhead.}
The gated policy's median per-query latency of 143.7\,s against the
closed-book 7.42\,s reflects the cost of the gating drafts: the ensemble issues
the entropy/margin draft plus the two hallucination-probe drafts before the
final answer, and on CPU these dominate the per-query time. This overhead is
reported honestly as the price of the current ensemble and is itself an
argument for the calibration work and for the low-tier probe-only
configuration, where fewer gate calls are made.

%  Required preamble packages (shared with week1_implementation.tex):
%    \usepackage{booktabs} \usepackage{amsmath} \usepackage{amssymb}
%
%  NOTE: the prototype results table (Table~\ref{tab:p5-prototype}) is a
%  PLACEHOLDER pending completion of the 200-question P5 run. Fill the marked
%  cells from results/raw_logs/p5_mirage200.jsonl once the run finishes.
% =====================================================================

\section{Adaptive Retrieval Gating (Week 2)}
\label{sec:impl-week2}

This section documents the project's core novel contribution: the
training-free retrieval gates that decide, per query, whether to retrieve
external evidence or answer from the model's parametric knowledge. Building on
the foundation layer of Section~\ref{sec:impl-week1}, this sprint implemented
three independent gating signals, combined them into a majority-vote ensemble,
and composed them into the selective-retrieval policy~(P5). The result is a
working prototype in which P5 runs end-to-end on the local quantised model and
is directly comparable to the closed-book floor established in Week~1.

% ---------------------------------------------------------------------
\subsection{Design Rationale: Why Three Gates}
\label{sec:gate-rationale}

A single uncertainty signal is fragile. Token entropy can be inflated by
benign lexical variation (synonyms, formatting tokens) on a query the model
actually knows; the logit margin can be large on a confidently \emph{wrong}
answer; self-reported confidence is poorly calibrated in small models. Each
signal therefore has a characteristic failure mode in which it would gate
incorrectly. The design response is an ensemble of three signals derived from
\emph{different} aspects of the model's behaviour, combined by majority vote so
that no single signal's failure mode dominates the decision. This mirrors the
project's safety-first framing: in a clinical setting, the cost of an
unnecessary retrieval (latency, energy) is far lower than the cost of a
confidently wrong closed-book answer, so the ensemble is deliberately biased
towards retrieving when the signals disagree.

% ---------------------------------------------------------------------
\subsection{Gate 1: Token Entropy}
\label{sec:gate-entropy}

The entropy gate generates a short draft answer (48 tokens) without retrieval
and measures the model's uncertainty from the full-vocabulary logit
distribution at each generated position. For generated token $t$ with softmax
distribution $p_t$ over the vocabulary $V$, the per-token Shannon entropy (in
nats) and the gate signal are
\begin{align}
  H(p_t) &= -\sum_{v \in V} p_t(v)\,\log p_t(v), &
  \bar{H} &= \frac{1}{N}\sum_{t=1}^{N} H(p_t),
\end{align}
and the gate retrieves when $\bar{H} > \tau_H$ (default $\tau_H = 2.5$). High
mean entropy indicates the model is spreading probability mass over many
candidate continuations --- a direct measure of lexical-level uncertainty.

Beyond the scalar $\bar{H}$ used for the decision, the gate also returns the
full per-token array $\{H(p_t)\}_{t=1}^N$. This array is the basis of the
\emph{Token-Level Entropy Attribution} contribution: it localises uncertainty
to specific tokens of the draft, revealing \emph{which part} of a multi-part
clinical question the model is unsure about rather than merely \emph{that} it
is unsure. The array is propagated through the policy result into the JSONL
audit log, where it is available to the planned interpretability
visualisation. The gate requires the raw logit buffer
($\texttt{logits\_all=True}$); when that buffer is unavailable on the low
hardware tier, the gate abstains rather than producing a spurious signal.

% ---------------------------------------------------------------------
\subsection{Gate 2: Logit Margin}
\label{sec:gate-margin}

The margin gate measures the model's \emph{decisiveness} from the gap between
its two most probable tokens at each position. For generated token $t$ with
top-two probabilities $p_t^{(1)} \ge p_t^{(2)}$,
\begin{align}
  m_t &= p_t^{(1)} - p_t^{(2)}, &
  \bar{M} &= \frac{1}{N}\sum_{t=1}^{N} m_t,
\end{align}
and the gate retrieves when $\bar{M} < \tau_M$ (default $\tau_M = 0.3$). A
small margin means the model is torn between two near-equiprobable
continuations, a complementary notion to entropy: whereas entropy summarises
the whole distribution, the margin focuses on the decision boundary between the
top two candidates.

Crucially, the margin gate obtains its signal from the completion API's
top-$k$ log-probabilities rather than the raw logit buffer, a lighter
dependency than the entropy gate's. The parser that converts the API's
per-token log-probability dictionaries into margins was validated against the
real inference library during this sprint, confirming that the live output
shape matches the mock used in testing, so that the unit tests are faithful to
production behaviour. Basis: the margin signal of TARG.

% ---------------------------------------------------------------------
\subsection{Gate 3: Hallucination Probe}
\label{sec:gate-probe}

The hallucination-probe gate replaces the verbalized-confidence gate of the
original design. Rather than asking the model to introspect on its own
confidence --- a task at which small models are unreliable --- it tests whether
the model's answer is \emph{stable} under reframing. Two drafts are generated
from the same question under two different prompt framings:
\begin{align*}
  \text{draft}_A &= \mathcal{M}\bigl(\text{probe}_A(q)\bigr), &
  \text{draft}_B &= \mathcal{M}\bigl(\text{probe}_B(q)\bigr),
\end{align*}
where $\text{probe}_A$ adopts an assistant-role framing and $\text{probe}_B$ a
direct-question framing. For multiple-choice questions the gate compares the
extracted option letters and retrieves when they differ; for open-ended
questions it computes the token-level $F_1$ between the two drafts and
retrieves when $F_1 < 0.7$. The intuition is that genuinely held knowledge is
invariant to surface framing, whereas a guessed answer is sensitive to it, so
disagreement across framings is evidence of instability.

This gate is purely text-based: it never reads the logit buffer, and is
therefore the only gate that remains available on the low hardware tier where
$\texttt{logits\_all=False}$. Its cost is two draft generations per query, an
overhead reported explicitly in the evaluation. Basis: self-consistency (Wang
et al., 2023), here reduced from majority voting over many samples to a binary
agreement test between two framings.

% ---------------------------------------------------------------------
\subsection{The Ensemble and Hardware-Aware Degradation}
\label{sec:gate-ensemble}

The three gates are combined by majority vote. Let $g_i \in
\{\textsc{retrieve}, \textsc{skip}\}$ be the decision of member $i$ among the
\emph{available} members $A \subseteq \{\text{entropy}, \text{margin},
\text{probe}\}$. The ensemble retrieves when
\[
  \bigl|\{\, i \in A : g_i = \textsc{retrieve} \,\}\bigr| \;\ge\; v_{\min},
  \qquad v_{\min} = 2 .
\]
Only members reporting themselves available cast a vote; a gate that abstains
(the entropy and margin gates on the low tier) is excluded from both the
numerator and the count of voters.

This availability accounting produces a principled \emph{degraded mode}. On the
low tier, only the probe gate is available, so the strict threshold
$v_{\min}=2$ is unreachable. Rather than failing, the ensemble falls back to
``retrieve if any available member votes to retrieve'' and records a
\texttt{degraded} flag in the audit log. The fallback is enforced at
construction time by the policy factory, which inspects the model's
$\texttt{logits\_all}$ setting and downgrades any requested gate to
probe-only when the logit buffer is disabled. The consequence is an important
design property: \emph{the gate a policy actually runs is constrained by the
hardware tier, not merely by the policy configuration file}. This makes the
hardware-aware behaviour auditable and turns the low-tier limitation into a
documented finding supporting the \emph{Hardware-Aware Policy Escalation}
contribution.

% ---------------------------------------------------------------------
\subsection{The Gated Policy (P5)}
\label{sec:p5-policy}

The selective-retrieval policy composes the ensemble with the retriever and the
language model. For each question it first obtains the gate decision; on
\textsc{retrieve} it behaves like the always-retrieve baseline (retrieve
evidence, answer with it in context), and on \textsc{skip} it behaves like the
closed-book baseline (answer from parametric knowledge alone). The complete
gate decision payload --- the scalar signal, the per-member votes, the
per-token entropy array, and the two probe drafts --- is attached to the policy
result and written into the per-query JSONL record, so that every gating
decision is fully reconstructable after the fact. Policy construction is
centralised in a factory that assembles the correct policy and gate from a
validated configuration, replacing the hand-wired construction used in Week~1.

% ---------------------------------------------------------------------
\subsection{Verification}
\label{sec:gate-verification}

Every gate and the gated policy were developed test-first against the mock
backend, whose synthetic peaked and uniform logit distributions exercise both
the \textsc{skip} and \textsc{retrieve} paths deterministically. The gate test
suite covers each signal's threshold behaviour, the per-token entropy
recording, the abstain paths under disabled logits, the ensemble vote
arithmetic, and the degraded-mode fallback; the policy test suite covers both
branches of P5, the propagation of gate fields into the result, and the
factory's construction of the correct gate per hardware tier (including the
low-tier probe-only downgrade). The full unit suite stands at 86 passing tests
and continues to execute without any model download, preserving the project's
fast, reproducible verification story.

% ---------------------------------------------------------------------
\subsection{Prototype Result}
\label{sec:p5-prototype}

The prototype was exercised by running P5 over the same 200 MIRAGE questions
used for the Week~1 closed-book baseline, using the local Llama-3.2-3B model
and a curated pilot retrieval corpus. Using the identical question set makes the
gated policy directly comparable to the closed-book floor of
Table~\ref{tab:p3-baseline}. The pilot corpus is a deliberately small,
curated set; retrieval-\emph{quality} figures obtained against it are
provisional and will be re-measured against the full corpus, but the
retrieval-\emph{rate}, gate-agreement and latency-overhead figures reported
here characterise the gating mechanism itself and are not sensitive to corpus
size. The aggregate prototype figures are given in
Table~\ref{tab:p5-prototype}, and the per-gate behaviour --- which proves the
most informative result of this sprint --- in Table~\ref{tab:p5-gatevotes}.

\begin{table}[ht]
  \centering
  \caption{Prototype P5 (3-gate ensemble, default thresholds) vs.\ the P3
           closed-book floor on the same 200 MIRAGE questions, Llama-3.2-3B
           Q4\_K\_M, medium tier. At the literature-default thresholds the gate
           never fired; the two policies are therefore identical in accuracy.
           See Table~\ref{tab:p5-gatevotes} for the cause.}
  \label{tab:p5-prototype}
  \begin{tabular}{lrr}
    \toprule
    \textbf{Metric} & \textbf{P3 (closed-book)} & \textbf{P5 (gated, uncal.)} \\
    \midrule
    Accuracy            & 62.0\% & 62.0\% \\
    Retrieval rate      & 0\%    & 0\% \\
    Median latency      & 7.42\,s & 143.7\,s \\
    95th-pct.\ latency  & 9.25\,s & 154.1\,s \\
    Mean energy / query & $6.77\times10^{-4}$\,kWh & $1.20\times10^{-3}$\,kWh \\
    \bottomrule
  \end{tabular}
\end{table}

\begin{table}[ht]
  \centering
  \caption{Per-gate signal statistics and vote counts over the 200 questions.
           The entropy and margin signals never cross their default thresholds
           on this quantised model, so only the hallucination probe ever votes
           to retrieve --- and a lone probe vote cannot reach the majority of
           two. This is a threshold-transfer finding, not a defect.}
  \label{tab:p5-gatevotes}
  \begin{tabular}{lrrrr}
    \toprule
    \textbf{Gate} & \textbf{Signal range} & \textbf{Default $\tau$} &
    \textbf{Retrieve votes} & \textbf{Skip votes} \\
    \midrule
    Entropy (nats)        & 0.27--1.59 & $>2.5$  & 0   & 200 \\
    Margin                & 0.45--0.89 & $<0.3$  & 0   & 200 \\
    Hallucination probe   & ---        & ---     & 80  & 120 \\
    \bottomrule
  \end{tabular}
\end{table}

\paragraph{Interpretation and the central finding.}
At the thresholds inherited from the general-domain literature ($\tau_H=2.5$
nats, $\tau_M=0.3$), the gate retrieved on \emph{zero} of the 200 questions, so
P5 reduced to the closed-book policy and reproduced its 62.0\% accuracy exactly.
The per-gate trace in Table~\ref{tab:p5-gatevotes} explains why: on the
4-bit-quantised Llama-3.2-3B, greedy drafts are sharply peaked, so the mean
token entropy (observed range 0.27--1.59 nats) never approaches the 2.5-nat
trigger, and the top-two probability margin (0.45--0.89) never falls below the
0.3 trigger. Only the hallucination probe --- which is threshold-free and
measures cross-framing consistency --- ever signalled uncertainty, voting to
retrieve on 80 of 200 questions; but under a majority rule requiring two of
three votes, a single dissenting gate can never carry the decision.

This is the most important result of the sprint, and it is a genuine empirical
finding rather than an implementation fault: \emph{gate thresholds calibrated
on general-domain models do not transfer to a small quantised medical model.}
It directly motivates the threshold-calibration sweep (Sprint~4): the observed
signal distributions suggest model-specific operating points near the
$75^{\text{th}}$ percentile of entropy ($\approx 0.9$ nats) and the
$25^{\text{th}}$ percentile of margin ($\approx 0.64$), which would move the
logit-based gates off their floor and let the three signals genuinely compete.
The prototype has thus done exactly what a prototype should: it validated the
end-to-end mechanism (gates compute, votes aggregate, the policy branches and
logs correctly) while surfacing a calibration requirement that becomes a
substantive point of analysis in the evaluation.

\paragraph{Latency overhead.}
The gated policy's median per-query latency of 143.7\,s against the
closed-book 7.42\,s reflects the cost of the gating drafts: the ensemble issues
the entropy/margin draft plus the two hallucination-probe drafts before the
final answer, and on CPU these dominate the per-query time. This overhead is
reported honestly as the price of the current ensemble and is itself an
argument for the calibration work and for the low-tier probe-only
configuration, where fewer gate calls are made.
